% GENERATED FILE. Edit canonical lessons, then run npm run generate:books.
% canonical-source-hash: fnv1a64:9c94b26e5a3243e6
% canonical-lessons: MR-C259-anolkhi, MR-C259-sadha, MR-C259-guntaguntica, MR-C259-avasyak, MR-C259-nirupyogi, MR-R259-first-pass-words-from-chapters-251-254, MR-R259-second-pass-words-from-chapters-255-259

\chapter{Unknown, Simple, Complicated, Necessary, Useless}
\label{ch:mr-unknown-simple-complicated-necessary-useless}
\hlchaptermodality{259}

\begin{chapteropening}
\textbf{By the end of this chapter:} \emph{I can say unknown, simple, complicated, necessary and useless.}

Complete the last lesson of chapter 259: I can say unknown, simple, complicated, necessary and useless.
\end{chapteropening}

\section[\texorpdfstring{\mr{अनोळखी} (anoḷkhī)}{अनोळखी (anoḷkhī)}]{\mr{अनोळखी} (anoḷkhī) — unknown, unfamiliar}

\label{lesson:MR-C259-anolkhi}



\begin{quote}
Before the new one: say the Marathi for all night, then the Marathi for all day.
\end{quote}

\subsection*{You'll want to know: \mr{अनोळखी}}
\textbf{\mr{अनोळखी}} — \emph{anoḷkhī} — \textquotedblleft{}unknown, unfamiliar\textquotedblright{}.

\begin{grammarlens}[title={What you've built}]
One more word for this chapter.
\end{grammarlens}

\subsection*{Your turn}
\begin{itemize}
\raggedright
  \item \emph{Say it:} \emph{anoḷkhī}
  \item \emph{Say it:} \emph{anoḷkhī}, once more
  \item \emph{Say it:} say \emph{anoḷkhī}
  \item \emph{Recall:} say the Marathi for all night, then the Marathi for all day, then say \emph{anoḷkhī} again
\end{itemize}

\begin{tcolorbox}[breakable,colback=teal!4,colframe=teal!35!black,title={Before you move on}]
What does \mr{अनोळखी} mean? (\textquotedblleft{}Unknown, unfamiliar\textquotedblright{}.) Say it once more.
\end{tcolorbox}
\section[\texorpdfstring{\mr{साधा} / \mr{साधी} (sādhā / sādhī)}{साधा / साधी (sādhā / sādhī)}]{\mr{साधा} / \mr{साधी} (sādhā / sādhī) — simple, plain}

\label{lesson:MR-C259-sadha}



\begin{quote}
Before the new one: say the Marathi for all day, then the Marathi for unknown.
\end{quote}

\subsection*{You'll want to know: \mr{साधा} / \mr{साधी}}
\textbf{\mr{साधा} / \mr{साधी}} — \emph{sādhā / sādhī} — \textquotedblleft{}simple, plain\textquotedblright{}.

\begin{grammarlens}[title={What you've built}]
One more word for this chapter.
\end{grammarlens}

\subsection*{Your turn}
\begin{itemize}
\raggedright
  \item \emph{Say it:} \emph{sādhā / sādhī}
  \item \emph{Say it:} \emph{sādhā / sādhī}, once more
  \item \emph{Say it:} say \emph{sādhā / sādhī}
  \item \emph{Recall:} say the Marathi for all day, then the Marathi for unknown, then say \emph{sādhā / sādhī} again
\end{itemize}

\begin{tcolorbox}[breakable,colback=teal!4,colframe=teal!35!black,title={Before you move on}]
What does \mr{साधा} / \mr{साधी} mean? (\textquotedblleft{}Simple, plain\textquotedblright{}.) Say it once more.
\end{tcolorbox}
\section[\texorpdfstring{\mr{गुंतागुंतीचा} (guntāguntīcā)}{गुंतागुंतीचा (guntāguntīcā)}]{\mr{गुंतागुंतीचा} (guntāguntīcā) — complicated}

\label{lesson:MR-C259-guntaguntica}



\begin{quote}
Before the new one: say the Marathi for unknown, then the Marathi for simple.
\end{quote}

\subsection*{You'll want to know: \mr{गुंतागुंतीचा}}
\textbf{\mr{गुंतागुंतीचा}} — \emph{guntāguntīcā} — \textquotedblleft{}complicated\textquotedblright{}.

\begin{grammarlens}[title={What you've built}]
One more word for this chapter.
\end{grammarlens}

\subsection*{Your turn}
\begin{itemize}
\raggedright
  \item \emph{Say it:} \emph{guntāguntīcā}
  \item \emph{Say it:} \emph{guntāguntīcā}, once more
  \item \emph{Say it:} say \emph{guntāguntīcā}
  \item \emph{Recall:} say the Marathi for unknown, then the Marathi for simple, then say \emph{guntāguntīcā} again
\end{itemize}

\begin{tcolorbox}[breakable,colback=teal!4,colframe=teal!35!black,title={Before you move on}]
What does \mr{गुंतागुंतीचा} mean? (\textquotedblleft{}Complicated\textquotedblright{}.) Say it once more.
\end{tcolorbox}
\section[\texorpdfstring{\mr{आवश्यक} (āvaśyak)}{आवश्यक (āvaśyak)}]{\mr{आवश्यक} (āvaśyak) — necessary}

\label{lesson:MR-C259-avasyak}



\begin{quote}
Before the new one: say the Marathi for simple, then the Marathi for complicated.
\end{quote}

\subsection*{You'll want to know: \mr{आवश्यक}}
\textbf{\mr{आवश्यक}} — \emph{āvaśyak} — \textquotedblleft{}necessary\textquotedblright{}.

\begin{grammarlens}[title={What you've built}]
One more word for this chapter.
\end{grammarlens}

\subsection*{Your turn}
\begin{itemize}
\raggedright
  \item \emph{Say it:} \emph{āvaśyak}
  \item \emph{Say it:} \emph{āvaśyak}, once more
  \item \emph{Say it:} say \emph{āvaśyak}
  \item \emph{Recall:} say the Marathi for simple, then the Marathi for complicated, then say \emph{āvaśyak} again
\end{itemize}

\begin{tcolorbox}[breakable,colback=teal!4,colframe=teal!35!black,title={Before you move on}]
What does \mr{आवश्यक} mean? (\textquotedblleft{}Necessary\textquotedblright{}.) Say it once more.
\end{tcolorbox}
\section[\texorpdfstring{\mr{निरुपयोगी} (nirupyogī)}{निरुपयोगी (nirupyogī)}]{\mr{निरुपयोगी} (nirupyogī) — useless}

\label{lesson:MR-C259-nirupyogi}



\begin{quote}
Before the new one: say the Marathi for complicated, then the Marathi for necessary.
\end{quote}

\subsection*{You'll want to know: \mr{निरुपयोगी}}
\textbf{\mr{निरुपयोगी}} — \emph{nirupyogī} — \textquotedblleft{}useless\textquotedblright{}.

\begin{grammarlens}[title={What you've built}]
One more word for this chapter.
\end{grammarlens}

\subsection*{Your turn}
\begin{itemize}
\raggedright
  \item \emph{Say it:} \emph{nirupyogī}
  \item \emph{Say it:} \emph{nirupyogī}, once more
  \item \emph{Say it:} say \emph{nirupyogī}
  \item \emph{Recall:} say the Marathi for complicated, then the Marathi for necessary, then say \emph{nirupyogī} again
\end{itemize}

\begin{tcolorbox}[breakable,colback=teal!4,colframe=teal!35!black,title={Before you move on}]
What does \mr{निरुपयोगी} mean? (\textquotedblleft{}Useless\textquotedblright{}.) Say it once more.
\end{tcolorbox}
\section[(dialogue)]{Words from chapters 251-254 — a first pass}

\label{lesson:MR-R259-first-pass-words-from-chapters-251-254}



\begin{quote}
Say the words for necessary and useless.
\end{quote}

\subsection*{Your turn}
Say these aloud:

\begin{itemize}
\raggedright
  \item to try, to keep going, to bring to a stop, to move, to heat
  \item to have fun, to get angry, to worry, to smoke, to go with
  \item to turn, to enjoy, to imagine, to inform, to complain
  \item to steal, to discover, to serve, to bear, to dream
\end{itemize}

\begin{tcolorbox}[breakable,colback=teal!4,colframe=teal!35!black,title={Before you move on}]
To try? (\textbf{\mr{प्रयत्न} \mr{करणे}}, \emph{prayatna karṇe}.) To keep going? (\textbf{\mr{चालू} \mr{ठेवणे}}, \emph{cālū ṭhevṇe}.) To bring to a stop? (\textbf{\mr{थांबवणे}}, \emph{thāmbavṇe}.) To move? (\textbf{\mr{हलवणे}}, \emph{halavṇe}.) To heat? (\textbf{\mr{गरम} \mr{करणे}}, \emph{garam karṇe}.) To have fun? (\textbf{\mr{मजा} \mr{करणे}}, \emph{majā karṇe}.) To get angry? (\textbf{\mr{रागावणे}}, \emph{rāgāvṇe}.) To worry? (\textbf{\mr{काळजी} \mr{करणे}}, \emph{kāḷjī karṇe}.) To smoke? (\textbf{\mr{धूम्रपान} \mr{करणे}}, \emph{dhūmrapān karṇe}.) To go with? (\textbf{\mr{सोबत} \mr{जाणे}}, \emph{sobat jāṇe}.) To turn? (\textbf{\mr{वळणे}}, \emph{vaḷṇe}.) To enjoy? (\textbf{\mr{आनंद} \mr{घेणे}}, \emph{ānand gheṇe}.) To imagine? (\textbf{\mr{कल्पना} \mr{करणे}}, \emph{kalpanā karṇe}.) To inform? (\textbf{\mr{कळवणे}}, \emph{kaḷavṇe}.) To complain? (\textbf{\mr{तक्रार} \mr{करणे}}, \emph{takrār karṇe}.) To steal? (\textbf{\mr{चोरणे}}, \emph{corṇe}.) To discover? (\textbf{\mr{शोधून} \mr{काढणे}}, \emph{śodhūn kāḍhṇe}.) To serve? (\textbf{\mr{सेवा} \mr{करणे}}, \emph{sevā karṇe}.) To bear? (\textbf{\mr{सहन} \mr{करणे}}, \emph{sahan karṇe}.) To dream? (\textbf{\mr{स्वप्न} \mr{पाहणे}}, \emph{svapna pāhṇe}.)
\end{tcolorbox}
\section[(dialogue)]{Words from chapters 255-259 — a second pass}

\label{lesson:MR-R259-second-pass-words-from-chapters-255-259}



\begin{quote}
Say the words for to insist and to interrupt.
\end{quote}

\subsection*{Your turn}
Say these aloud:

\begin{itemize}
\raggedright
  \item to insist, to interrupt, to mix, to pray, to recommend
  \item to borrow, to look after, to express, to get thin, to feed
  \item rest, anxiety, friendship, trust, courage
  \item nowadays, a century, an era, all night, all day
  \item unknown, simple, complicated, necessary, useless
\end{itemize}

\begin{tcolorbox}[breakable,colback=teal!4,colframe=teal!35!black,title={Before you move on}]
To insist? (\textbf{\mr{आग्रह} \mr{करणे}}, \emph{āgrah karṇe}.) To interrupt? (\textbf{\mr{मध्ये} \mr{बोलणे}}, \emph{madhye bolṇe}.) To mix? (\textbf{\mr{मिसळणे}}, \emph{misaḷṇe}.) To pray? (\textbf{\mr{प्रार्थना} \mr{करणे}}, \emph{prārthanā karṇe}.) To recommend? (\textbf{\mr{शिफारस} \mr{करणे}}, \emph{śiphāras karṇe}.) To borrow? (\textbf{\mr{उसने} \mr{घेणे}}, \emph{usne gheṇe}.) To look after? (\textbf{\mr{जपणे}}, \emph{japṇe}.) To express? (\textbf{\mr{व्यक्त} \mr{करणे}}, \emph{vyakta karṇe}.) To get thin? (\textbf{\mr{बारीक} \mr{होणे}}, \emph{bārīk hoṇe}.) To feed? (\textbf{\mr{खाऊ} \mr{घालणे}}, \emph{khāū ghālṇe}.) Rest? (\textbf{\mr{विश्रांती}}, \emph{viśrāntī}.) Anxiety? (\textbf{\mr{चिंता}}, \emph{cintā}.) Friendship? (\textbf{\mr{मैत्री}}, \emph{maitrī}.) Trust? (\textbf{\mr{विश्वास}}, \emph{viśvās}.) Courage? (\textbf{\mr{धैर्य}}, \emph{dhairya}.) Nowadays? (\textbf{\mr{हल्ली}}, \emph{hallī}.) A century? (\textbf{\mr{शतक}}, \emph{śatak}.) An era? (\textbf{\mr{युग}}, \emph{yug}.) All night? (\textbf{\mr{रात्रभर}}, \emph{rātrabhar}.) All day? (\textbf{\mr{दिवसभर}}, \emph{divasbhar}.) Unknown? (\textbf{\mr{अनोळखी}}, \emph{anoḷkhī}.) Simple? (\textbf{\mr{साधा} / \mr{साधी}}, \emph{sādhā / sādhī}.) Complicated? (\textbf{\mr{गुंतागुंतीचा}}, \emph{guntāguntīcā}.) Necessary? (\textbf{\mr{आवश्यक}}, \emph{āvaśyak}.) Useless? (\textbf{\mr{निरुपयोगी}}, \emph{nirupyogī}.)
\end{tcolorbox}
